\section*{実験に⽤いたソースコード}

解析に用いたスクリプトのうち，本研究の手法を特徴づける部分を抜粋する．
全体は \url{https://github.com/KurenNagata/bravo} に置いてある．

\subsection*{A. 実プレー時間の復元（\texttt{extract\_dataset.py}）}

イベントに記録された継続時間をつなぎ，10秒を超える空白を停止とみなして時計から除く．
これにより，前後の「3分」を停止時間の混入なくそろえられる．

\begin{verbatim}
MAX_LIVE_GAP_SEC = 10.0

def live_segments(events, period):
    """possession 内のイベント列から実プレー区間を復元する。"""
    by_possession = defaultdict(list)
    for e in events:
        if e["period"] != period or "possession" not in e:
            continue
        if e["type"]["name"] in NON_PLAY_TYPES:
            continue
        s = float(tsec(e))
        duration = max(0.0, float(e.get("duration", 0.0) or 0.0))
        by_possession[e["possession"]].append((s, s + duration))

    pieces = []
    for intervals in by_possession.values():
        intervals.sort()
        cur_s, cur_t = intervals[0]
        for s, t in intervals[1:]:
            if s - cur_t > MAX_LIVE_GAP_SEC:   # 10秒超の空白は停止
                if cur_t > cur_s:
                    pieces.append((cur_s, cur_t))
                cur_s, cur_t = s, t
            else:
                cur_t = max(cur_t, t)
        if cur_t > cur_s:
            pieces.append((cur_s, cur_t))

    pieces.sort()
    merged = []
    for s, t in pieces:
        if not merged or s > merged[-1][1]:
            merged.append((s, t))
        else:
            merged[-1] = (merged[-1][0], max(merged[-1][1], t))
    return merged
\end{verbatim}

\subsection*{B. 32項目の組み立て（\texttt{train\_models.py}）}

イベントの記録から機械的に列挙した項目を，内容・位置，実プレー1分あたりの回数，
テンポの3種類に分けてまとめる．指標を人が選び直さないための処理である．

\begin{verbatim}
def feature_sets(all_cols):
    per_active_min = [c for c in all_cols
                      if c.endswith("_per_active_min")]
    raw_counts = [c for c in all_cols
                  if (c.startswith("n_") or c == "xg_sum")
                  and not c.endswith("_per_active_min")]
    tempo = [
        "events_per_possession",
        "passes_per_possession",
        "within_possession_gap_mean",
    ]
    content = [c for c in all_cols
               if c not in per_active_min
               and c not in raw_counts
               and c not in tempo]
    return {
        "flow_full": content + per_active_min + tempo,  # 32項目
        "content_only": content,
    }
\end{verbatim}

\subsection*{C. fold先行の対応づけ（\texttt{did\_foldfirst.py}）}

試合を先に5つの組へ分け，対応づけを同じ組の中だけで行う．
これにより，同じ試合・同じ対が学習と評価にまたがらないことを設計だけで保証する．

\begin{verbatim}
def match_within_blocks(treated, controls, block_of, use_xg):
    used, pairs = set(), []
    order = sorted(range(len(treated)),
                   key=lambda i: treated[i]["clock"])
    for i in order:
        t = treated[i]
        if use_xg and t.get("xg") is None:
            continue
        tb = block_of[str(t["match_id"])]
        best_j, best_cost = None, None
        for j, c in enumerate(controls):
            if j in used or block_of[str(c["match_id"])] != tb:
                continue                    # 同じ組の中だけで探す
            if c["period"] != t["period"]:
                continue                    # 前後半を一致
            if c["state"] != t["state"]:
                continue                    # 点差状態を一致
            if c["is_home"] != t["is_home"]:
                continue                    # ホーム／アウェーを一致
            if c["match_id"] == t["match_id"]:
                continue
            gap = abs(c["clock"] - t["clock"])
            if gap > MATCH_MINUTE_TOL:       # 時刻差 8 分以内
                continue
            cost = gap / MATCH_MINUTE_TOL
            if use_xg:
                if c.get("xg") is None:
                    continue
                dxg = abs(c["xg"] - t["xg"])
                if dxg > XG_TOL:             # xG 差 0.10 以内
                    continue
                cost += dxg / XG_TOL
            if best_cost is None or cost < best_cost:
                best_j, best_cost = j, cost
        if best_j is not None:
            used.add(best_j)
            pairs.append((t, controls[best_j]))
    return pairs
\end{verbatim}

\subsection*{D. 差の差の学習・評価（\texttt{did\_foldfirst.py}，\texttt{train\_models.py}）}

後区間と前区間の差を入力とし，組をそのままfoldとして交差検証する．
AUCは対になった処置側・対照側の判定値の大小から求め，
95\%信頼区間は連結成分を単位とするブートストラップで与える．

\begin{verbatim}
# 判別器（train_models.py）
Pipeline([
    ("impute", SimpleImputer(strategy="median")),
    ("scale", StandardScaler()),
    ("clf", LogisticRegression(max_iter=2000, C=1.0,
                               random_state=SEED)),
])

# 学習・評価（did_foldfirst.py の evaluate_blocks より抜粋）
rows = []
for t, c in pairs:
    block = block_of[str(t["match_id"])]
    for label, u in ((1, t), (0, c)):
        delta = {f: (u["after"][f] - u["before"][f])   # 後 - 前
                 for f in cols}
        rows.append({"block": block, "pair": t["key"],
                     "label": label, **delta})
df = pd.DataFrame(rows)

probs = np.empty(len(df))
for b in range(N_BLOCKS):              # 組をそのまま fold にする
    te = df["block"] == b
    model = build_models()[PRIMARY_MODEL].fit(
        df.loc[~te, cols].to_numpy(dtype=float),
        df.loc[~te, "label"].to_numpy())
    probs[te.to_numpy()] = model.predict_proba(
        df.loc[te, cols].to_numpy(dtype=float))[:, 1]

piv = df.assign(prob=probs).pivot_table(
    index="pair", columns="label", values="prob").dropna()
values = piv[[0, 1]].to_numpy()
observed = roc_auc_score(np.tile([0, 1], len(values)),
                         values.reshape(-1))

# 連結成分を単位としたブートストラップ 5,000 回
rng = np.random.default_rng(SEED)
boots = np.empty(N_BOOT)
for b in range(N_BOOT):
    pick = rng.choice(keys, size=len(keys), replace=True)
    idx = [i for cmp_ in pick for i in by_comp[cmp_]]
    boots[b] = auc_of(values[idx])
ci95 = np.quantile(boots, [0.025, 0.975])
\end{verbatim}
